
%(BEGIN_QUESTION)
% Copyright 2011, Tony R. Kuphaldt, released under the Creative Commons Attribution License (v 1.0)
% This means you may do almost anything with this work of mine, so long as you give me proper credit

Les og lag en disposisjon av Case History \#71 (``The Amazing Problem Free Plant'') fra Michael Browns samling av veiledninger om optimalisering av reguleringssløyfer. Forbered deg på en gjennomtenkt diskusjon med lærer og medstudenter om begrepene og eksemplene i denne lesningen, og svar på følgende spørsmål:

\begin{itemize}
\item{} I starten av denne rapporten viser Mr. Brown til et anslag der et anlegg med 30\% av sine automatiske reguleringer som går ineffektivt ville få en {\it alternativkostnad} på \$350,000 årlig. Forklar hva uttrykket ``alternativkostnad'' betyr.
\vskip 10pt
\item{} Identifiser trekk i grafene i figur 1 som viser {\it integral} regulatorvirkning samt {\it proporsjonal} regulatorvirkning. Forklar også hvordan vi kan avgjøre fra denne grafen at reguleringsventilen har problemer, og forklar hvordan vi kan bestemme regulatorens virkemåte (direkte eller revers) fra trendene.
\vskip 10pt
\item{} Undersøk trenden i figur 3, og identifiser problemet som vises der. Forklar hvordan vi kan se at dette er et ``metningsproblem'' i reguleringsventilen og ikke et ``metningsproblem'' i transmitteren.
\vskip 10pt
\item{} På sløyfe \#7 i dette anlegget -- en nivåreguleringssløyfe -- anbefalte Mr. Brown å bruke filtrering på PV-signalet for å håndtere stor mengde støy. Hvorfor ikke bare redusere proporsjonalvirkningen i regulatoren og heller bruke aggressiv integralvirkning?
\end{itemize}

\vskip 20pt \vbox{\hrule \hbox{\strut \vrule{} {\bf Forslag til sokratisk diskusjon} \vrule} \hrule}

\begin{itemize}
\item{} ``Alternativkostnad'' er et svært nyttig økonomisk begrep, med mange praktiske anvendelser både i industri og i hverdagslivet (personlig økonomi). Identifiser noen alternativkostnader i egen erfaring. Hint: hva er alternativkostnaden(e) ved å gå på skole på heltid?
\item{} Trenden i figur 1 avslører ganske mye om regulatorens PID-innstilling. Undersøk PV-, SP- og utgangstrendene og forklar hva vi kan utlede om regulatorinnstillingen. Er den P-, I- eller D-dominert? Kan vi beregne regulatorens forsterkning, reset-tid eller rate-tid?
\item{} I Mr. Browns analyse av sløyfe \#5 konkluderer han med: ``Either the valve was passing, or it was stroked incorrectly.'' Forklar hva denne uttalelsen betyr.
\end{itemize}

\underbar{file i01550}
%(END_QUESTION)





%(BEGIN_ANSWER)


%(END_ANSWER)





%(BEGIN_NOTES)

{\it Alternativkostnad} er et økonomisk begrep som betyr verdien av det alternativet jeg gir avkall på for å gjøre det jeg gjør nå.

\vskip 10pt

Figur 1: proporsjonalvirkning kan sees ganske tydelig i responsen på SP-endringen. Integralrespons forklarer den rampende utgangen mens PV $>$ SP. Like før neste SP-endring ser du at PV faktisk kommer lik SP.

\vskip 10pt

Figur 3: PV responderer ikke når ventilposisjonen overstiger omtrent 80\%. Siden PV fortsatt er innenfor måleområdet, vet vi at dette sannsynligvis ikke er et transmitterproblem, men mer sannsynlig et ventilmetningsproblem (eller eventuelt et problem med ventilposisjoneren).

\vskip 10pt

På sløyfe \#7 (nivåregulering med støyete PV) ble aggressiv proporsjonalvirkning kombinert med PV-filtrering funnet å være den beste løsningen. Å redusere regulatorforsterkningen og la integral gjøre det meste av arbeidet ville gjort at sløyfen kunne fungere med liten eller ingen filtrering, men det ville gitt oversving etter SP-endringer, slik det alltid blir med integralvirkning i integrerende sløyfer.

 


\vskip 20pt \vbox{\hrule \hbox{\strut \vrule{} {\bf Forslag til sokratisk diskusjon} \vrule} \hrule}

\begin{itemize}
\item{} Hvorfor foreslo Michael Brown å begrense regulatorens utgang til 80\% på sløyfe \#6 (figur 3)?
\item{} Bortsett fra metningen vist i figur 3, viser reguleringsventilen i denne sløyfen andre betydelige problemer?
\end{itemize}

%INDEX% Reading assignment: Michael Brown Case History #71, "The amazing problem free plant"

%(END_NOTES)

